% ===================================================================
%  ТАБЛИЦА № 6 — Дом внутри. — Мебель. — Еда. — Освещение.
%  Traduction 2026 d'apres texto/io, controlee sur texto/fr, texto/en
%  et texto/es. Consigne : voir texto/ru/10-tablica-01.tex.
%
%  Deux notes, toutes deux marquees « (*) », et deux appels : celui de
%  « babo » au bloc c1-12-1, celui de « lambrequino » au bloc c2-01-1.
%  L'ordre les departage, comme dans les autres colonnes.
%
%  LA FEMME DE CHAMBRE. L'ido du bloc c1-06-1 porte « (6) » la ou le
%  francais porte « (16) », et c'est le francais qui a raison. Le
%  russe suit le francais et ecrit (16), comme l'anglais et
%  l'espagnol ; gravuri/korekti.json redresse le renvoi ido dans la
%  page de lecture, et le gros plan s'ouvre dans les sept colonnes.
% ===================================================================

%%K t06-apar-1 apar
\VUsaut{6.0mm}
\VUtitre{15.0pt}{60}{ТАБЛИЦА № 6}
\VUsaut{1.4mm}
\VUtitre{11.4pt}{0}{Дом внутри, Мебель, Еда,}
\VUsaut{0.4mm}
\VUtitre{11.4pt}{0}{Освещение.}
\VUsaut{2.2mm}

%%K t06-apar-2 apar
Дом готов. Дядя Ивана переехал в своё новое жилище, расположение которого
мы уже знаем. Посмотрим теперь, как он его обставил. Сперва мы посетим:

%%K t06-tit-1 sub
\VUpk{10.2pt}{40}{Спальню.}
\VUsaut{3.0mm}

%%K t06-c1-01-1 p
1. --- Мы пришли в неудачную минуту, потому что служанка убирает
\VUgras{комнату}.

%%K t06-c1-02-1 p
2. --- На \VUgras{умывальнике} \textsuperscript{(1)} лежат там и сям разные вещи,
которыми моются, причёсываются и вообще приводят себя в порядок. Это
\VUgras{таз} \textsuperscript{(2)} и \VUgras{кувшин} \textsuperscript{(3)},
\VUgras{щётка для волос} \textsuperscript{(4)}, \VUgras{гребень} \textsuperscript{(5)},
\VUgras{мыло} \textsuperscript{(6)} и \VUgras{губка} \textsuperscript{(7)}, наконец
\VUgras{пузырёк} \textsuperscript{(8)} с полосканьем для рта.

%%K t06-c1-03-1 p
3. --- На полу, перед умывальником, стоит \VUgras{ведро} \textsuperscript{(9)}
для грязной воды.

%%K t06-c1-04-1 p
4. --- В \VUgras{прихожей} \textsuperscript{(10)} видны несколько
\VUgras{вешалок} \textsuperscript{(13)} и два \VUgras{зонта} \textsuperscript{(11)} в
\VUgras{подставке} \textsuperscript{(12)}.

%%K t06-c1-05-1 p
5. --- По другую сторону двери стоит \VUgras{зеркальный шкаф}
\textsuperscript{(14)}, в котором хранится \VUgras{бельё} \textsuperscript{(15)}.

%%K t06-c1-06-1 p
6. --- Воспользуемся тем, что \VUgras{горничная} \textsuperscript{(16)} стелет
\VUgras{постель} \textsuperscript{(17)}, и рассмотрим её части. \VUgras{Кровать}
\textsuperscript{(20)} несёт добротный \VUgras{пружинный матрац} \textsuperscript{(21)},
на который Жюстина сейчас положит шерстяной \VUgras{тюфяк} \textsuperscript{(22)}.
Потом она возьмёт со стульев две \VUgras{простыни} \textsuperscript{(23)} и
\VUgras{одеяло} \textsuperscript{(24)}. Затем положит в изголовье \VUgras{валик}
\textsuperscript{(25)} и \VUgras{подушку} \textsuperscript{(26)}, которые лежат вместе с
\VUgras{периной} \textsuperscript{(27)} на \VUgras{прикроватном коврике}
\textsuperscript{(32)}. Метёлкой она обмахнёт \VUgras{занавеси} \textsuperscript{(19)} и
\VUgras{полог} \textsuperscript{(18)}, и часть работы сделана.

%%K t06-c1-07-1 p
7. --- Потом ей придётся немного прибрать \VUgras{ночной столик}
\textsuperscript{(28)}, завести \VUgras{будильник} \textsuperscript{(29)}, приготовить
\VUgras{ночник} \textsuperscript{(30)} и унести \VUgras{свечу} \textsuperscript{(31)},
чтобы вычистить подсвечник.

%%K t06-tit-2 sub
\VUsaut{5.0mm}
\VUpk{10.2pt}{40}{Рабочая корзинка и каминный прибор.}
\VUsaut{3.0mm}

%%K t06-c1-08-1 p
8. --- \VUgras{Рабочая корзинка} \textsuperscript{(34)} возле \VUgras{табурета}
\textsuperscript{(33)} тоже сильно нуждается в уборке. Придётся сложить
\VUgras{вышивку} \textsuperscript{(35)}, убрать в \VUgras{рабочий столик}
\textsuperscript{(38)} \VUgras{напёрсток}, \VUgras{нитки} \textsuperscript{(36)},
\VUgras{иголки} \textsuperscript{(37)} и \VUgras{ножницы} \textsuperscript{(39)} и
запереть крышку столика \VUgras{ключом} \textsuperscript{(40)}.

%%K t06-c1-09-1 p
9. --- На \VUgras{круглом столике} \textsuperscript{(41)} посреди комнаты Жюстина
сменит воду в \VUgras{графине} \textsuperscript{(42)}. Она положит \VUgras{сахар}
в \VUgras{сахарницу} \textsuperscript{(43)}, вычистит \VUgras{щипчики для сахара}
\textsuperscript{(44)} и унесёт \VUgras{платяную щётку} \textsuperscript{(46)}.

%%K t06-c1-10-1 p
10. --- Правая сторона комнаты даст ей меньше работы, потому что
\VUgras{кушетка} \textsuperscript{(47)} и её \VUgras{подушка} \textsuperscript{(48)}
на своих местах, а так как не холодно, ничего не трогали и в приборе
\VUgras{камина} \textsuperscript{(49)}. \VUgras{Совок} \textsuperscript{(50)},
\VUgras{щипцы} \textsuperscript{(51)}, \VUgras{таганы} \textsuperscript{(55)} и
\VUgras{решётка для золы} \textsuperscript{(56)} чисты; \VUgras{каминный экран}
\textsuperscript{(54)} не сдвинут, а \VUgras{ящик для дров} \textsuperscript{(52)}
полон \VUgras{дров} \textsuperscript{(53)}.

%%K t06-noto-1 noto
\VUnotes{143.2mm}{%
(*) Babo = маленький ребёнок; англ. \textit{baby}, фр. \textit{bébé},
ит. \textit{bambino}.%
}

%%K t06-noto-2 noto
\VUnotes{143.2mm}{%
(*) Lambrequino = фр. \textit{lambrequin}, исп.
\textit{lambrequines}, порт. \textit{lambrequins}, и даже нем. и англ.
\textit{lambrequins} --- то есть и по-немецки, и по-английски, и
по-французски, и по-португальски, и по-испански.%
}

%%K t06-c1-11-1 p
11. --- Ей ещё надо будет обмахнуть \VUgras{каминные часы} \textsuperscript{(57)},
\VUgras{канделябры} \textsuperscript{(58)} и \VUgras{подносики для мелочей}
\textsuperscript{(59)}, а потом протереть \VUgras{зеркало} \textsuperscript{(60)}.

%%K t06-c1-12-1 p
12. --- Что до \VUgras{колыбели} \textsuperscript{(61)} малыша (бабо (*)), она
уже готова.

%%K t06-tit-3 sub
\VUsaut{5.0mm}
\VUpk{10.2pt}{40}{Гостиную.}
\VUsaut{3.0mm}

%%K t06-c2-01-1 p
1. --- А вот и гостиная. Это очень красивая \VUgras{комната}. Камин в
правом углу украшен изящной \VUgras{статуэткой} \textsuperscript{(1)} и двумя
красивыми \VUgras{вазами} \textsuperscript{(3)}, а задрапирован бархатным
\VUgras{ламбрекеном} \textsuperscript{(2)} (*).

%%K t06-c2-02-1 p
2. --- Чтобы искры из очага не подожгли ковёр, перед камином поставлен
медный \VUgras{экран} \textsuperscript{(4)}.

%%K t06-c2-03-1 p
3. --- Рядом стоит раскрытое изящное дамское \VUgras{бюро} \textsuperscript{(5)}.
Здесь \VUgras{хозяйка дома} \textsuperscript{(23)} пишет свои письма.

%%K t06-c2-04-1 p
4. --- Окно убрано \VUgras{тюлевыми занавесками} \textsuperscript{(6)},
подхваченными \VUgras{подхватами} \textsuperscript{(8)} на \VUgras{крючках}
\textsuperscript{(7)}, и тяжёлыми матерчатыми шторами с \VUgras{подзором}
\textsuperscript{(9)} сверху.

%%K t06-c2-05-1 p
5. --- В нише окна \VUgras{жардиньерка} \textsuperscript{(11)} держит зелёное
\VUgras{растение} \textsuperscript{(10)}, великолепную пальму, листья которой
достают почти до \VUgras{керосиновой лампы} \textsuperscript{(12)}.

%%K t06-c2-06-1 p
6. --- \VUgras{Абажур} \textsuperscript{(13)} сделала Мария, прелестная
\VUgras{барышня} \textsuperscript{(14)} за \VUgras{роялем} \textsuperscript{(15)}.
Сидя очень прямо на \VUgras{табурете} \textsuperscript{(16)}, она разучивает
пьесу. Она очень искусна и очень аккуратна, о чём говорит её
\VUgras{нотный шкафчик} \textsuperscript{(17)}: там полный порядок.

%%K t06-tit-4 sub
\VUsaut{5.0mm}
\VUpk{10.2pt}{40}{Шалуна.}
\VUsaut{3.0mm}

%%K t06-c2-07-1 p
7. --- Её брат Павел ищет \VUgras{книгу} \textsuperscript{(19)} в
\VUgras{книжном шкафу} \textsuperscript{(18)}. Это \VUgras{мальчик}
\textsuperscript{(20)} непоседливый и совсем несносный. Повиснув на шкафу, чтобы
достать слишком высоко стоящую книгу, он рискует свалить стоящий наверху
\VUgras{бюст} \textsuperscript{(21)}.

%%K t06-c2-08-1 p
8. --- Он пользуется для этой глупости тем, что мать занята разговором с
приятельницей, \VUgras{дамой} \textsuperscript{(22)}, приехавшей погостить у неё
несколько дней. Мать сидит на \VUgras{диване} \textsuperscript{(24)} возле
\VUgras{кресла} \textsuperscript{(25)}, а её приятельница на \VUgras{стуле}
\textsuperscript{(27)}, от которого нам видна только \VUgras{спинка}
\textsuperscript{(26)}.

%%K t06-c2-09-1 p
9. --- В большой \VUgras{раме} \textsuperscript{(30)} на стене --- \VUgras{портрет}
\textsuperscript{(29)} родственника. В \VUgras{альбоме} \textsuperscript{(32)},
лежащем на столе, собраны \VUgras{фотографии} \textsuperscript{(33)} остальной
родни. Дети только что его листали, потому что \VUgras{стулья}
\textsuperscript{(31)} ещё сдвинуты вокруг стола.

%%K t06-c2-10-1 p
10. --- Фарфоровую \VUgras{китайскую вазу} \textsuperscript{(34)} возле
\VUgras{ножа для бумаги} \textsuperscript{(35)} подарили Марии на именины, а
\VUgras{ширму} \textsuperscript{(36)}, занимающую угол комнаты, привёз из Китая
её дядя.

%%K t06-c2-11-1 p
11. --- Оживлённая красивыми \VUgras{картинами} \textsuperscript{(37)} на
\VUgras{стене} \textsuperscript{(42)}, освещённая потолочной \VUgras{люстрой}
\textsuperscript{(38)}, чьи \VUgras{электрические лампочки} \textsuperscript{(39)} так
весело светят по вечерам, эта гостиная кажется очень приятной. Мягкий
\VUgras{ковёр} \textsuperscript{(40)}, разостланный по полу, и такой удобный
\VUgras{электрический звонок} \textsuperscript{(41)} делают её вполне уютной.

%%K t06-tit-5 sub
\VUsaut{3.6mm}
\VUpk{10.2pt}{40}{Столовую.}
\VUsaut{3.0mm}

%%K t06-c3-01-1 p
1. --- Прозвонили к обеду. Вся семья, кроме Марии, собралась вокруг
стола. Столовая приятно натоплена отличной изразцовой \VUgras{печью}
\textsuperscript{(1)} с тонкой \VUgras{трубой} \textsuperscript{(2)}.

%%K t06-c3-02-1 p
2. --- В этой комнате немного мебели. Вдоль стены, над \VUgras{табуретом}
\textsuperscript{(4)}, висит нарядный \VUgras{умывальный фонтанчик}
\textsuperscript{(3)}. Рядом стоит \VUgras{буфет} \textsuperscript{(5)}, на который
ставят холодные блюда, десерты и разную посуду, пока не нужную.

%%K t06-c3-03-1 p
3. --- Так, на верхней полке мы видим \VUgras{жареную курицу}
\textsuperscript{(7)}, \VUgras{рыбу} \textsuperscript{(8)} и \VUgras{вазу}
\textsuperscript{(10)} с \VUgras{фруктами} \textsuperscript{(9)}. Ниже стоят
\VUgras{сыр} \textsuperscript{(11)}, \VUgras{хлеб} \textsuperscript{(12)} и
\VUgras{судок} \textsuperscript{(13)}, в котором есть всё, что нужно, чтобы
заправить салат: масло, уксус, \VUgras{соль} \textsuperscript{(14)} и
\VUgras{перец} \textsuperscript{(15)}. Наконец, на нижней полке стоит
\VUgras{пирог} \textsuperscript{(17)} рядом с \VUgras{горчичницей}
\textsuperscript{(16)}.

%%K t06-c3-04-1 p
4. --- Часы в столовой, которые называются \VUgras{стенными}
\textsuperscript{(20)}, очень необычной формы. Они вставлены в деревянную раму,
в которой есть также \VUgras{барометр} \textsuperscript{(19)} и
\VUgras{термометр} \textsuperscript{(18)}.

%%K t06-tit-6 sub
\VUsaut{4.6mm}
\VUpk{10.2pt}{40}{Как подают обед.}
\VUsaut{3.0mm}

%%K t06-c3-05-1 p
5. --- Кругом всей комнаты стены обшиты \VUgras{деревянными панелями}
\textsuperscript{(21)} из вощёного дуба. Матерчатая \VUgras{портьера}
\textsuperscript{(22)} довольно хорошо закрывает дверь, ведущую на веранду.
Туда семья пойдёт после обеда пить кофе. \VUgras{Чашки} \textsuperscript{(24)}
и \VUgras{блюдца} \textsuperscript{(25)} уже расставлены на \VUgras{горке}
\textsuperscript{(23)}.

%%K t06-c3-06-1 p
6. --- Над столом к потолку прикреплена \VUgras{висячая лампа}
\textsuperscript{(26)}; её очень яркий свет заливает по вечерам всю столовую.

%%K t06-c3-07-1 p
7. --- Посмотрим теперь на сам \VUgras{обеденный стол} \textsuperscript{(27)}.
Когда семья вошла, стол был накрыт. На \VUgras{скатерти} \textsuperscript{(28)}
поставили на каждое место \VUgras{тарелку} \textsuperscript{(33)},
\VUgras{салфетку} \textsuperscript{(29)}, \VUgras{ложку} \textsuperscript{(34)},
\VUgras{вилку} \textsuperscript{(35)}, \VUgras{нож} \textsuperscript{(36)} и
\VUgras{стакан} \textsuperscript{(32)}. Были там и \VUgras{бутылки}
\textsuperscript{(30)} для вина и \VUgras{графин} \textsuperscript{(31)} для воды.

%%K t06-c3-08-1 p
8. --- \VUgras{Слуга} \textsuperscript{(39)} сперва внёс \VUgras{суповую миску}
\textsuperscript{(37)}, в которой ещё дымится суп. Теперь он подаёт первое
\VUgras{блюдо} \textsuperscript{(38)}, мясное. Потом он обнесёт всех теми, что мы
уже назвали; и наконец, после \VUgras{десерта}, воспользуется тем, что
хозяева перешли на веранду, чтобы убрать со стола и снова привести
столовую в порядок.

%%K t06-c3-08-2 p
\textit{Примечание.} --- \textit{Обиходные слова о еде даны здесь только
в общих чертах. Список будет дополнен в двух таблицах «Гостиница»
(№ 13) и «Рынок» (№ 15).}

%%K t06-tit-7 sub
\VUsaut{4.6mm}
\VUpk{10.2pt}{40}{Кухню.}
\VUsaut{3.0mm}

%%K t06-c4-01-1 p
1. --- На кухне, куда мы заглядываем через приотворённую дверь, нас
встречает живописный беспорядок. Перед \VUgras{очагом} \textsuperscript{(1)}, где
трещит \VUgras{огонь} \textsuperscript{(3)}, установлен \VUgras{вертельный станок}
\textsuperscript{(5)}. Прекрасное \VUgras{жаркое} \textsuperscript{(7)} надето на
\VUgras{вертел} \textsuperscript{(6)} и доходит.

%%K t06-c4-02-1 p
2. --- \VUgras{Мехи} \textsuperscript{(8)} и \VUgras{совок для угля}
\textsuperscript{(9)}, которыми только что пользовались, остались на полу вместе
с \VUgras{поленьями} \textsuperscript{(10)}.

%%K t06-c4-03-1 p
3. --- На каминной полке порядок: \VUgras{фонарь} \textsuperscript{(11)},
\VUgras{спичечница} \textsuperscript{(13)} со \VUgras{спичками} \textsuperscript{(12)},
\VUgras{подсвечник} \textsuperscript{(14)} и \VUgras{ночной подсвечник}
\textsuperscript{(15)} со своей \VUgras{свечой} \textsuperscript{(16)} --- всё на своих
местах. Зато большое \VUgras{ведро для угля} \textsuperscript{(18)}, полное
\VUgras{угля} \textsuperscript{(17)}, за которым \VUgras{кухарка} \textsuperscript{(23)}
только что спускалась в подвал, оставлено посреди кухни.

%%K t06-c4-04-1 p
4. --- Дело в том, что бедной женщине надо торопиться, чтобы завтрак
поспел вовремя. Сейчас она у \VUgras{плиты} \textsuperscript{(19)}, помешивает
\VUgras{соус} \textsuperscript{(24)}, которому нельзя дать пригореть.

%%K t06-c4-05-1 p
5. --- В \VUgras{духовке} \textsuperscript{(20)} печётся пирог, который она
сделала детям к полднику. Над плитой висят \VUgras{половник}
\textsuperscript{(21)} и \VUgras{шумовка} \textsuperscript{(22)}.

%%K t06-tit-8 sub
\VUsaut{4.0mm}
\VUpk{10.2pt}{40}{Кухонную посуду.}
\VUsaut{3.0mm}

%%K t06-c4-06-1 p
6. --- \VUgras{Набор кастрюль} \textsuperscript{(25)}, висящий в глубине комнаты,
очень чист. Красивые медные \VUgras{кастрюли} \textsuperscript{(26)},
\VUgras{молочник} \textsuperscript{(27)}, \VUgras{дуршлаг} \textsuperscript{(28)} и
\VUgras{тёрка} \textsuperscript{(29)} тщательно вычищены.

%%K t06-c4-07-1 p
7. --- \VUgras{Солонка} \textsuperscript{(30)} стоит на маленьком поставце рядом
с \VUgras{чайником} \textsuperscript{(31)} и \VUgras{бутылью для масла}
\textsuperscript{(32)}. В \VUgras{ящике} \textsuperscript{(33)}, надо думать, лежат
приборы и кухонные ножи.

%%K t06-c4-08-1 p
8. --- \VUgras{Метла} \textsuperscript{(34)} и \VUgras{метёлка из перьев}
\textsuperscript{(35)} стоят в углу. Кухарка только что пользовалась
\VUgras{колодой} \textsuperscript{(36)}; она оставила её у окна.

%%K t06-c4-09-1 p
9. --- \VUgras{Полотенце} \textsuperscript{(37)} висит на стене рядом с
\VUgras{воронкой} \textsuperscript{(38)}. В этой части кухни держат
\VUgras{решётку} \textsuperscript{(39)}, \VUgras{сечку} \textsuperscript{(40)} и
\VUgras{разделочную доску} \textsuperscript{(41)}. Затем, над сечкой, на
\VUgras{полке} \textsuperscript{(42)}, стоят \VUgras{горшки} \textsuperscript{(43)},
\VUgras{чугунок} \textsuperscript{(44)} со своей \VUgras{крышкой}
\textsuperscript{(45)} и \VUgras{котёл} \textsuperscript{(46)}. \VUgras{Сковорода}
\textsuperscript{(47)} висит рядом.

%%K t06-c4-10-1 p
10. --- Только медный \VUgras{кран} \textsuperscript{(48)} бросает отблеск в этот
угол. \VUgras{Раковина} \textsuperscript{(49)}, на которой стоит большой
\VUgras{кувшин} \textsuperscript{(50)}, должно быть очень удобна со своим
шкафчиком, где держат \VUgras{бочонок с уксусом} \textsuperscript{(51)} с его
\VUgras{краником} \textsuperscript{(52)}, \VUgras{ящик для мусора}
\textsuperscript{(53)} и \VUgras{грязную посуду} \textsuperscript{(54)}.

%%K t06-tit-9 sub
\VUsaut{4.0mm}
\VUpk{10.2pt}{40}{Что ещё держат на кухне.}
\VUsaut{3.0mm}

%%K t06-c4-11-1 p
11. --- Большая \VUgras{лохань} \textsuperscript{(55)}, в которой моют
\VUgras{овощи} \textsuperscript{(58)}, и чугунный \VUgras{котёл} \textsuperscript{(56)},
стоящий на \VUgras{изразцовом полу} \textsuperscript{(57)}, тоже, верно,
принадлежат этому шкафчику.

%%K t06-c4-12-1 p
12. --- В плитах у кухарки, конечно, недостатка нет, потому что вот тут,
на углу стола, стоит \VUgras{газовая горелка} \textsuperscript{(59)}, а на ней
греется вода в кастрюльке. \VUgras{Пар} \textsuperscript{(60)}, выходящий из неё,
говорит нам, что вода, должно быть, кипит. Эта вода, верно, для кофе,
потому что вот на другом конце стола стоят \VUgras{кофейник}
\textsuperscript{(64)} и \VUgras{кофейная мельница} \textsuperscript{(63)}.

%%K t06-c4-13-1 p
13. --- Горелка соединена с \VUgras{газовым рожком} \textsuperscript{(62)}
длинной \VUgras{резиновой трубкой} \textsuperscript{(61)}.

%%K t06-c4-14-1 p
14. --- \VUgras{Подёнщица} \textsuperscript{(66)}, приходящая помогать кухарке,
готовит овощи к обеду. Когда они будут очищены и вымыты, она положит их
в \VUgras{кастрюлю} \textsuperscript{(65)} до той поры, пока не придёт время их
варить.

%%K t06-tit-10 sub
\VUsaut{4.0mm}
{\centering\fontsize{10.2pt}{10.2pt}\selectfont\textls[40]{\scshape Ванную.} \textit{(Комната для купанья.)}\par}
\VUsaut{3.0mm}

%%K t06-c5-01-1 p
1. --- Остаётся заглянуть ещё в одно место, чтобы узнать главные комнаты
дома: посмотреть на устройство ванной. Это недолго, потому что она
невелика, и мы быстро увидим, что у \VUgras{ванны} \textsuperscript{(1)} есть
хорошая \VUgras{колонка} \textsuperscript{(2)}, что \VUgras{мыльница}
\textsuperscript{(3)} у купающегося под рукой, а \VUgras{вешалка для полотенец}
\textsuperscript{(5)} должна облегчать сушку \VUgras{полотенец} \textsuperscript{(4)}.

%%K t06-c5-02-1 p
2. --- У этой комнаты стены покрыты \VUgras{глазурованной плиткой}
\textsuperscript{(6)}, что позволило поставить \VUgras{душ} \textsuperscript{(7)}, не
опасаясь, что брызжущая при пользовании вода их испортит. Маленький
цинковый \VUgras{таз} \textsuperscript{(8)} для ножных ванн довершает обстановку.
\VUsaut{6.0mm}
\VUfilet{22mm}
